% Companion to report.tex: every equation explained for oral presentation.
\documentclass[11pt,a4paper]{article}
\usepackage[T1]{fontenc}
\usepackage{newtxtext,newtxmath}
\usepackage[a4paper,margin=0.9in]{geometry}
\usepackage{setspace}\setstretch{1.12}
\usepackage{amsmath,booktabs,tabularx,enumitem,xcolor,microtype}
\usepackage[most]{tcolorbox}
\usepackage[hidelinks]{hyperref}
\setlength{\parindent}{0pt}\setlength{\parskip}{5pt}
\newcommand{\code}[1]{{\small\texttt{#1}}}
\newtcolorbox{say}{colback=black!4,colframe=black!45,boxrule=0.4pt,arc=1.5mm,left=4pt,right=4pt,top=2pt,bottom=2pt,
	title={\small\bfseries How to say it in 20 seconds},coltitle=black,colbacktitle=black!12,fontupper=\small}
\newtcolorbox{ex}{colback=blue!3,colframe=blue!40,boxrule=0.4pt,arc=1.5mm,left=4pt,right=4pt,top=2pt,bottom=2pt,
	title={\small\bfseries Worked example (numbers checked by script)},coltitle=black,colbacktitle=blue!10,fontupper=\small}
\usepackage{titlesec}
\titleformat{\section}{\large\bfseries}{}{0pt}{}
\titlespacing*{\section}{0pt}{12pt}{4pt}

\begin{document}
\begin{center}
	{\LARGE\bfseries Equation Guide for the Sentinel Report}\\[4pt]
	{\large Md.\ Shifat Hasan, Roll 2107067 \,|\, CSE 4102}
\end{center}

Each section matches one equation of Chapter~III (numbers in brackets are the report's equation numbers). For every equation: what it does, what each symbol means, why it has that shape, a numeric example, the exact place in the code, and a short script for the viva. All formulas were re-checked against the source files; four statements in the first draft were corrected (the IK angle, the Cook--Torrance form, the normal-map scale and the wind formula) and the report now matches the code.

\section{G1. Noise deformer (3.1): turning a sphere into a rock}
\[\mathbf p'=\mathbf p+A\,\nu(f\,\mathbf p)\,\hat n,\qquad \lVert\mathbf p'-\mathbf p\rVert\le A.\]
\textbf{Meaning.} Every vertex $\mathbf p$ of a mesh is pushed along its surface direction $\hat n$ by a smooth random amount. \textbf{Symbols:} $A$ amplitude in metres (largest possible push); $\nu\in[-1,1]$ a 3-D fractal noise value; $f$ frequency, in cycles per metre (higher $f$ gives smaller bumps, because the noise is sampled at $f\mathbf p$); $\hat n$ unit normal of the original surface. \textbf{Why the bound holds:} $|\nu|\le1$ and $\lVert\hat n\rVert=1$, so the push is at most $A$; this is what makes object size ranges testable. Neighbouring vertices get nearly equal $\nu$ (noise is smooth), so the surface stays connected and lumpy, not spiky.
\begin{ex}$A=0.3$\,m and $\nu=0.5$ at one vertex: the vertex moves $0.3\times0.5=0.15$\,m outward. At $\nu=-1$ it moves 0.3\,m inward.\end{ex}
\textbf{Code:} \code{Engine/Procedural/Deform.odin}, case \code{Noise\_Displace}.
\begin{say}``A noise deformer moves each vertex along its normal by amplitude times smooth noise, so shapes get natural lumps while the maximum change is guaranteed to be the amplitude.''\end{say}

\section{G2. Two-bone inverse kinematics (3.2): placing a foot}
\[\cos\alpha=\frac{a^{2}+c^{2}-b^{2}}{2ac}\]
\textbf{Meaning.} A leg has an upper bone (length $a$, thigh) and a lower bone (length $b$, shin). Given where the foot must be, we need the hip angle. The hip, knee and foot form a triangle with sides $a$, $b$ and $c=|{\rm target}-{\rm hip}|$; the \emph{law of cosines} gives any angle from the three sides. \textbf{Symbols:} $\alpha$ angle at the hip between the hip-to-target line and the thigh. The knee position then lies at distance $a\cos\alpha$ along the target line and $a\sin\alpha$ to the side. The interior knee angle $\theta$ uses the same law: $\cos\theta=(a^2+b^2-c^2)/2ab$. \textbf{Safety:} the code clamps $c$ and the cosine to $[-1,1]$, so an unreachable target gives a straight leg instead of a numerical error.
\begin{ex}$a=b=0.45$\,m, foot target $c=0.8$\,m away: $\cos\alpha=0.64/0.72=0.889$, so $\alpha=27.3^{\circ}$ and the knee angle is $125.5^{\circ}$ (bent but not folded). At $c=0.9=a+b$ the leg is straight ($\alpha=0$).\end{ex}
\textbf{Code:} \code{Game/Characters/Ik.odin} (law of cosines in the header comment and body).
\begin{say}``Hip, knee and foot make a triangle whose sides I know, so the law of cosines gives the joint angle; no animation files are needed.''\end{say}

\section{G3. Fractal noise (3.3) and why textures tile}
\[F(\mathbf x)=\frac{\sum_{i=0}^{n-1}2^{-i}\,\nu(2^{i}\mathbf x)}{\sum_{i=0}^{n-1}2^{-i}}\]
\textbf{Base noise $\nu$.} At each grid corner a random unit direction (one of 16, picked by a hash of the corner index) is stored. For a point inside a cell, $\nu$ takes the dot product of each corner direction with the vector from that corner to the point, then blends the four values with the smooth curve $s(t)=6t^5-15t^4+10t^3$ (it has zero slope and zero curvature at 0 and 1, so no grid lines show), and multiplies by $\sqrt2$ so the result spans $[-1,1]$.
\textbf{Fractal sum.} Octave $i$ has twice the frequency ($2^i\mathbf x$) and half the strength ($2^{-i}$): big hills, then boulders, then pebbles. The denominator is the total strength, so $F$ stays in $[-1,1]$.
\textbf{Tiling.} The corner direction is chosen from the \emph{cell index modulo the period} $P$ (and $P$ doubles each octave). Then $F(\mathbf x+P\mathbf e)=F(\mathbf x)$ for whole-number shifts $\mathbf e$: the right edge of a texture equals its left edge, so repeated tiles have no seam. The test suite checks this exactly (shift $(1,-2)$ tiles; difference must be 0).
\begin{ex}$n=4$ octaves: weights $1,\tfrac12,\tfrac14,\tfrac18$ sum to $1.875$. A value $\nu=+1$ in every octave gives $F=1$; typical values mix and give about $\pm0.4$. $s(0.5)=0.5$ exactly, and $s(0.25)=0.104$, so the blend is gentle near cell edges.\end{ex}
\textbf{Code:} \code{Shaders/Include/Noise.glsl} (\code{gradient\_noise}, \code{fbm}).
\begin{say}``Fractal noise adds smooth random layers, each twice as detailed and half as strong; by wrapping the grid index it repeats exactly, so textures tile.''\end{say}
\textbf{Variants used:} \emph{domain warping} $F(\mathbf x+s[F_1(\mathbf x),F_2(\mathbf x)])$ evaluates $F$ at a point shifted by two other noises, which bends straight lattice patterns into swirls (still periodic, since every term is). \emph{Ridged} noise uses $(1-|\nu|)^2$: $|\nu|$ folds the noise at zero so zero crossings become sharp ridges (cracks, bark). \emph{Cellular (Worley)} noise uses the distance to the nearest random point (leaves, pebbles).

\section{G4. Normal map from a height map (3.4)}
\[h_x\approx\frac{3(h_{NE}+h_{SE}-h_{NW}-h_{SW})+10(h_E-h_W)}{32\,\Delta\,N},\qquad \hat n=\frac{(-k\,h_x,\,-k\,h_y,\,1)}{\lVert(-k\,h_x,\,-k\,h_y,\,1)\rVert}\]
\textbf{Meaning.} A recipe provides a height $h$ per texel. Lighting only needs the surface \emph{direction}, which tilts away from the slope: a surface rising to the right ($h_x>0$) tilts its normal to the left ($-h_x$). \textbf{Scharr filter:} the slope is a weighted difference of the right column ($3,10,3$) minus the left column of the $3\times3$ neighbourhood. Weights sum to 16 over a distance of two texels, hence the divisor $32\Delta$; it treats diagonal and axis-aligned slopes equally, which the plain two-point difference does not. \textbf{Symbols:} $\Delta$ one texel (here $1/1024$ of the tile); $N$ cells per tile (so a bumpy recipe with more cells is not made steeper); $k$ the per-material bump strength; the vector is normalised to unit length. The last component $1$ means a flat area gives $\hat n=(0,0,1)$.
\begin{ex}A perfect ramp $h=0.7x$ (with $\Delta=1$): the numerator is $3\cdot2.8+10\cdot1.4=22.4$, divided by 32 gives $0.7$, the true slope (verified by script). With $k=1$: $\hat n=(-0.7,0,1)/1.22=(-0.57,0,0.82)$, tilted $35^{\circ}$.\end{ex}
\textbf{Cavity darkening.} Where $h$ is below the mean of its four neighbours (a negative Laplacian), the albedo is multiplied by $1-0.5c$ and the ambient occlusion by $1-c$, where $c\in[0,0.5]$ grows with the depth of the dent. \textbf{Triplanar weights:} $w_i=|n_i|^4/\sum_j|n_j|^4$ (the code uses exponent 4): a surface facing $+y$ mostly uses the top-down projection. For $\mathbf n=(0.6,0.8,0)$ the weights are $0.24,0.76,0$.
\textbf{Code:} \code{Shaders/Include/BakeMain.glsl}, \code{Shaders/Include/Triplanar.glsl}.
\begin{say}``The height map's slope tells how the surface tilts; a Scharr filter measures it fairly in every direction, and the normal is minus the slope with a one up, normalised.''\end{say}

\section{G5. Cook--Torrance reflection (3.5)}
\[f=\frac{(1-F_l)(1-F_v)\,\mathbf c_d}{\pi}+D\,V\,F,\quad D=\frac{\alpha^{2}}{\pi\big((\mathbf n\!\cdot\!\mathbf h)^{2}(\alpha^{2}-1)+1\big)^{2}},\quad F=F_0+(1-F_0)(1-\mathbf v\!\cdot\!\mathbf h)^{5}\]
\[V=\frac{0.5}{(\mathbf n\!\cdot\!\mathbf l)\sqrt{(\mathbf n\!\cdot\!\mathbf v)^2(1-\alpha^2)+\alpha^2}+(\mathbf n\!\cdot\!\mathbf v)\sqrt{(\mathbf n\!\cdot\!\mathbf l)^2(1-\alpha^2)+\alpha^2}}\]
\textbf{Meaning.} $f$ is the fraction of light arriving from direction $\mathbf l$ that leaves toward the viewer $\mathbf v$; the lighting pass multiplies it by the light's brightness and $\mathbf n\!\cdot\!\mathbf l$. The surface is imagined as many tiny mirrors (microfacets). \textbf{Vectors:} $\mathbf n$ surface normal, $\mathbf l$ to the light, $\mathbf v$ to the eye, $\mathbf h$ halfway between $\mathbf l$ and $\mathbf v$ (only facets facing $\mathbf h$ reflect light from $\mathbf l$ to $\mathbf v$).
\textbf{Specular part $DVF$.} $D$ (GGX): the share of facets facing $\mathbf h$; rough surfaces have small peak, wide lobe. $\alpha=\text{roughness}^2$. $V=G/(4\,\mathbf n\!\cdot\!\mathbf l\,\mathbf n\!\cdot\!\mathbf v)$: the geometry term $G$ (facets blocking each other) with the denominator folded in (height-correlated Smith form). $F$ (Schlick): every surface reflects more at glancing angles; $F_0$ is reflectance looking straight on ($0.04$ for plastic, stone, paint; the albedo colour for metals).
\textbf{Diffuse part.} $\mathbf c_d/\pi$ is the Lambert term ($\pi$ makes total reflected light equal the colour). It is scaled by $(1-F_l)(1-F_v)$ because light that mirror-reflected at the surface cannot also scatter inside; this keeps energy $\le1$. Metals have $\mathbf c_d=0$.
\begin{ex}Roughness 0.5 gives $\alpha=0.25$. Peak of $D$ at $\mathbf n\!\cdot\!\mathbf h=1$: $1/(\pi\cdot0.0625)=5.09$. Fresnel with $F_0=0.04$: looking at $60^{\circ}$ gives $F=0.07$; at $85^{\circ}$ (grazing) $F=0.65$, why wet roads shine at low angles.\end{ex}
\textbf{Code:} \code{Shaders/Include/Brdf.glsl}, \code{microfacet\_cook\_torrance}.
\begin{say}``Light reflects as a diffuse part plus a specular part from microscopic mirrors: GGX says how many face the eye, Smith says how many hide each other, Schlick says how reflective at this angle.''\end{say}

\section{G6. Fog by distance (haze)}
\[C=T\,C_{\text{surface}}+(1-T)\,C_{\text{haze}},\qquad T=e^{-kd},\quad k=3.5\times10^{-4}\ \text{m}^{-1}\]
\textbf{Meaning.} Air absorbs and scatters light. $T$ is the fraction of a surface's light that survives $d$ metres of air (Beer--Lambert law: equal thickness removes an equal \emph{fraction}, so the decay is exponential); the missing part is replaced by the sky-coloured haze, a linear mix. \textbf{Symbols:} $d$ distance to the camera; $k$ density.
\begin{ex}$d=100$\,m: $T=0.966$ (almost clear). $d=1$\,km: $T=0.705$. $d=3$\,km: $T=0.35$, so far hills are mostly sky colour, giving depth. The haze colour uses only the smooth atmosphere (never stars or the moon), which fixed the night streak bug.\end{ex}
\textbf{Code:} \code{Shaders/DeferredBase.glsl} (\code{FOG\_DENSITY\_PER\_METER}).
\begin{say}``Each metre of air removes the same percentage of light, so transmittance is an exponential, and what is lost is replaced by sky colour.''\end{say}

\section{G7. ACES tone mapping (3.6)}
\[C_{\text{out}}=\operatorname{clamp}_{[0,1]}\!\left(\frac{x(2.51x+0.03)}{x(2.43x+0.59)+0.14}\right),\qquad x=\text{exposure}\times\text{radiance}\]
\textbf{Meaning.} Scene light (HDR) can be 0.01 in shade or 40 at the sun, but a display shows 0 to 1. This curve (Narkowicz's fit of the ACES film curve) compresses highlights smoothly (shoulder), lifts dark detail gently (toe) and has an S-shaped contrast like film. \textbf{Symbols:} $x$ radiance scaled by exposure (1 by day, rising to 6 at night, like the eye adapting); applied to each colour channel. The numbers are fitted constants.
\begin{ex}$x=0.18$ (mid-grey) $\to0.267$; $x=1\to0.804$; $x=4\to0.973$. Quadrupling brightness from 1 to 4 changes the output by only 0.17: highlights do not clip to white abruptly.\end{ex}
After the curve the image is graded, converted to sRGB for the display and dithered by about one 8-bit step to avoid banding. \textbf{Code:} \code{Shaders/Include/Tonemap.glsl}, \code{Shaders/PostTonemap.glsl}.
\begin{say}``A rational curve maps unlimited brightness into zero-to-one with a soft shoulder, the way film does.''\end{say}

\section{G8. Sun height (inline in Section 3.2.5)}
\[\sin(\text{elev})=\sin\varphi\sin\delta+\cos\varphi\cos\delta\cos h,\qquad h=15^{\circ}(\text{hours}-12)\]
\textbf{Meaning.} Standard spherical astronomy. $\varphi=35^{\circ}$ latitude, $\delta=15^{\circ}$ the sun's declination (a late-spring day), $h$ the hour angle (the Earth turns $15^{\circ}$ per hour; $h=0$ at noon). The full direction vector uses the same terms with east and south components. Positive result: sun above the horizon.
\begin{ex}Noon: $\sin=0.1485+0.7913=0.9398$, elevation $70^{\circ}$. At 6:00 and 18:00: $8.5^{\circ}$ (low, orange sun). Midnight: $-40^{\circ}$, so the moon (opposite the sun) lights the scene.\end{ex}
\textbf{Code:} \code{Game/Gameplay/TimeOfDay.odin}.

\section{G9. Cloud plane and wind (3.7)}
\[\mathbf P_{xz}=\frac{H}{d_y}\,(d_x,d_z)+\mathbf w\,t\]
\textbf{Meaning.} Clouds are drawn as a flat noise layer at height $H$. For a view direction $\mathbf d=(d_x,d_y,d_z)$, a ray from the ground reaches height $H$ after travelling $H/d_y$ times the direction (similar triangles); that gives the horizontal position $(d_x,d_z)H/d_y$ where the noise is looked up. Adding $\mathbf wt$ ($\mathbf w$ wind velocity in m/s, $t$ time) slides the pattern, so clouds drift. A faster detail layer ($\times1.35$) and a cirrus layer at 7\,km ($\times1.8$ speed) make shapes change and give depth. Rays near the horizon ($d_y<0.04$) are skipped, because $H/d_y$ would blow up.
\begin{ex}$\mathbf d=(0.8,0.5,0.33)$ and $H=1500$: $H/d_y=3000$, point $(2400,995)$\,m. With $\mathbf w=(14,5)$\,m/s the pattern moved $(840,300)$\,m after one minute.\end{ex}
\textbf{Code:} \code{Shaders/Include/Sky.glsl} (\code{add\_clouds}, \code{cloud\_density}); wind in \code{Game/Sandbox/Sandbox.odin}.

\section{G10. Tree sway (cantilever wind)}
\[\Delta_x=2s\,(0.04+0.1\,g)\left(\frac{y}{9}\right)^{2},\quad \Delta_z=0.35\,\Delta_x,\quad g=\tfrac12+\tfrac12\sin(1.3\,t+0.05\,x+0.03\,z)\]
\textbf{Meaning.} The vertex shader shifts each tree vertex sideways. A real tree is like a beam fixed at the root: it barely moves at the bottom and most at the top, so the shift grows with the \emph{square} of height $y$ (clamped to 9\,m). $g\in[0,1]$ is a gust that travels across the land (position terms $0.05x+0.03z$ delay it from tree to tree, so trees do not sway in lockstep). $s$ is the per-mesh sway (0.6 for trunks, 1 for foliage; foliage above 1 also adds a small fast flutter). The factor 2 and offsets $0.04$/$0.1$ are tuned amplitudes: calm sway 0.08\,m, full gust 0.28\,m at the top.
\begin{ex}Foliage ($s=1$), full gust $g=1$: at $y=9$\,m, $\Delta_x=2(0.14)(1)=0.28$\,m; at half height $4.5$\,m it is $0.07$\,m (a quarter, because of the square).\end{ex}
\textbf{Code:} \code{Shaders/Geometry.glsl}, \code{wind\_offset}.

\section{G11. Vehicle steering (bicycle model) and engine power}
\[\dot\psi=\frac{v\tan\delta}{L},\qquad \dot v=a_0\left(1-0.6\,\frac{|v|}{v_{\max}}\right)\]
\textbf{Steering.} Treat the car as a bicycle with wheel base $L$ (m), front wheel angle $\delta$, speed $v$. The car turns on a circle of radius $R=L/\tan\delta$, so the yaw rate is $v/R=v\tan\delta/L$ (rad/s). Faster means turning faster; reversing flips the sign, as in a real car. (Tracked vehicles instead turn at a fixed rate and can pivot.)
\textbf{Power fade.} Engine thrust is power divided by speed, so acceleration is high from rest and falls as speed rises. $a_0$ is the maximum acceleration, $v_{\max}$ the limit; at top speed acceleration is still $0.4a_0$ but the speed is clamped to the limit.
\begin{ex}Jeep ($L=3$, $\delta_{\max}=0.55$\,rad, $v_{\max}=18$, $a_0=2.6$): at 10\,m/s full lock gives $\dot\psi=10\tan0.55/3=2.04$\,rad/s. Solving the speed equation, $v(t)=30(1-e^{-0.0867t})$: 2.5\,m/s after 1\,s, and 18\,m/s after 10.6\,s (it was 24\,m/s with 7\,m/s$^2$ before the realism change).\end{ex}
\textbf{Code:} \code{Engine/World/GroundVehicle.odin} (\code{yaw\_rate\_for}, \code{next\_speed}).

\section{Short answers to likely questions}
\begin{itemize}[leftmargin=0.25in,itemsep=1pt]
	\item \textbf{Why not use image textures?} The project rule forbids files; noise gives infinite variety from a seed and tiles exactly.
	\item \textbf{Why a quintic fade and not linear blending?} Linear blending shows grid lines (slope jumps at cell edges); $6t^5-15t^4+10t^3$ has zero first and second derivative there.
	\item \textbf{Why $\sqrt2$ in the noise?} Unit gradients bound the raw value by $\sqrt2/2\approx0.707$; multiplying by $\sqrt2$ makes the range $[-1,1]$.
	\item \textbf{Why exponential fog?} Beer--Lambert: each metre removes the same fraction of light.
	\item \textbf{Why is the GPU time 10.9\,ms but the frame 16.7\,ms?} The display waits for the 60\,Hz refresh; the extra 5.8\,ms is idle headroom.
	\item \textbf{What did you correct after checking?} IK used the root angle (law of cosines on $a,c,b$), Cook--Torrance has a visibility term $V$ and a diffuse part, the normal-map slope is divided by cells per tile, and the wind has exact constants.
\end{itemize}
\end{document}
